\documentclass[10pt,a4paper]{amsart}
\usepackage[margin=19mm]{geometry}
\usepackage{amsmath,amssymb,mathtools}
\usepackage[scr=rsfs,cal=euler]{mathalfa}
\usepackage{xfrac}
\usepackage[unicode,hidelinks,bookmarks=false]{hyperref}
\newcommand{\ie}{i.e.\ }
\newcommand{\cf}{cf.\ }
\newcommand{\resp}{respectively\ }
\newcommand{\Subsection}{\subsection}
\providecommand{\nobreakdash}{\nobreak\hbox{-}}
\setlength{\emergencystretch}{2em}
\multlinegap0pt
\usepackage{bm}
\usepackage{bbm}
\usepackage{dsfont}
\usepackage{xspace}
\usepackage{cancel}
\usepackage{mathtools}
\usepackage[shortlabels]{enumitem}
\usepackage{amsthm}
\makeatletter
\@ifundefined{thm}{\newtheorem{thm}{Thm}}{}
\@ifundefined{lem}{\newtheorem{lem}[thm]{Lemma}}{}
\@ifundefined{prop}{\newtheorem{prop}[thm]{Proposition}}{}
\makeatother
\title[BMT Independence]{BMT Independence}
\author{Octavio Arizmendi, Saul Rogelio Mendoza, Josu\'e Vazquez-Becerra}
\date{Source: arXiv:2309.04123v1 (CC BY 4.0)}
\begin{document}
\maketitle

\noindent\emph{Source and licence.} BMT Independence. Version arXiv:2309.04123v1 (CC BY 4.0), distributed under the Creative Commons Attribution 4.0 International licence, as recorded in the arXiv licence metadata for this exact version and shown on the abstract page https://arxiv.org/abs/2309.04123v1. Full rights notice: licenses/arxiv-2309.04123-CC-BY-4.0.txt.\par\smallskip

\noindent\textit{Editorial scope.} Counted unit: the Prop whose source label is \texttt{prop\_associativity} in the article (source0.tex lines 1074--1093, source label \texttt{prop\_associativity}), delivered together with the article's own complete original proof of it (source0.tex lines 1094--1282, one \texttt{proof} environment of 189 lines). No mathematics is written, repaired or re-derived: statement and proof are the article's own text line for line. Required context retained uncounted: lem\_weak\_bm (lines 1012--1021); lem\_weak\_bm\_ver\_kernel (lines 1034--1037). Every definition, display and label of those lines is retained unchanged. Omitted from this package and not redistributed: 1--1011, 1022--1033, 1038--1073, 1283--2724. Nothing inside a retained excerpt is omitted or altered: every retained body is delivered exactly as the article's lines are written, so no omission receipt of any kind accompanies this package. Citations: the retained excerpts carry no citation command at all, so no bibliography entry of the article is transcribed and no reference is redistributed. Reference adaptation: none was needed and none was made; every reference inside the delivered unit resolves inside it, so no unresolved reference or citation is produced. Printed slips kept literal: no printed word, symbol, hypothesis or number of the counted statement or of its proof was altered. Layout and class: the article's own class is \texttt{extarticle} with packages including inputenc, fontenc, amsmath, amsfonts, amssymb, amsthm, graphicx, cancel, bm, enumerate, xcolor, stmaryrd, hyperref, tikz; the wrapper is the lane's pinned \texttt{amsart} editorial wrapper at 10pt on A4, which restates the theorem-like environments the retained text needs and adds the article's own macro definitions that the retained text needs (none), with extra wrapper packages (bm, bbm, dsfont, xspace, cancel, mathtools, enumitem). The delivered numbering is the unit's own. Assets: no retained span contains a figure environment, a graphics-inclusion command, a PDF-page inclusion, a table or a TikZ picture; no image file is copied into this package, the package declares no asset dependency and no image file is redistributed with it. Licence: version arXiv:2309.04123v1 of this article is distributed under the Creative Commons Attribution 4.0 International licence, as recorded in the arXiv licence metadata for this exact version and shown on the abstract page https://arxiv.org/abs/2309.04123v1. A full rights notice is delivered at licenses/arxiv-2309.04123-CC-BY-4.0.txt. Characters: every retained excerpt is pure ASCII, so the delivered engine, XeLaTeX with its default Latin Modern text font, reports no missing character.
\begin{lem}\label{lem_weak_bm}
	%Let $(\mathcal{A},\varphi)$ be a non-commutative probability space. 
	%
	Suppose $(\mathcal{A}_i)_{i \in I}$ is a family of sub-algebras BMT independent with respect to a digraph $G = (I,E)$. 
	%
	Let $a_1 a_2 \cdots a_m$ be a product (not necessarily alternating) of elements of  $(\mathcal{A}_i)_{i \in I}$ with $a_k \in \mathcal{A}_{i_k}$ and $\bm{i} = (i_1, i_2,\ldots,i_m)$. 
	%
	If there exist $1 \leq \ell < \ell' \leq m$ so that $\ker_{G}[\bm{i}|_S]$ and $\ker_{G}[\bm{i}|_{S^c}]$ are contained in $\ker_G[\bm{i}]$ with $S = \{ \ell,\ell+1,\cdots,\ell'-1\}$, then 
	$$\varphi(a_1  \cdots a_m) = \varphi(a_{\ell} \cdots a_{\ell'-1}) \varphi(a_1\cdots a_{\ell-1} a_{\ell'} \cdots a_m)$$
\end{lem}

\begin{lem}\label{lem_weak_bm_ver_kernel}
	Let $G = (I,E)$ be a digraph. For any tuple $\bm{i} = (i_1,i_2,\ldots,i_m)$ with $i_k \in I$ and any subset $S \subset [m]$, we have 
	$\ker_{G}[\bm{i}|_S] \subset \ker_G[\bm{i}]$ if and only if $k_1 \sim_{ \ker_G[\bm{i}]} k_2$ for any $k_1,k_2 \in S$ with $k_1 \sim_{ \ker_G[\bm{i}|_S]} k_2$ and $k_1 \nsim_{ \ker_G[\bm{i}]} k_2$ for all $k_1 \in S$ and $k_2 \in [m]\setminus S$. 
\end{lem}

\begin{prop}\label{prop_associativity}
	Let $(\mathcal{A},\varphi)$ be a non-commutative probability space. 
	%
	Suppose $(\mathcal{A}_i)_{i \in I}$ is a family of subalgebras bmt independent with respect to a digraph $G = (I,E)$. 
	%
	If $\{ I_j : j \in J \}$ is a partition of $I$ into non-empty pairwise disjoint subsets and $\mathcal{B}_j = \mathrm{alg}( \mathcal{A}_i : i \in I_j )$, then  
	\begin{enumerate}[(i)]
		
		
		\item $(\mathcal{B}_j)_{j \in J}$ are tensor independent if $(i,i') \in E$ whenever $i \in I_{j}$ and $i' \in I_{j'}$ where $j,j' \in J$ with $j \neq j'$  
		
		\item $(\mathcal{B}_j)_{j \in J}$ are boolean independent if $(i,i') \notin E$ whenever $i \in I_{j}$ and $i' \in I_{j'}$ where $j,j' \in J$ with $j \neq j'$  
		
		
		\item $(\mathcal{B}_j)_{j \in J}$ are monotone independent if $J$ has a total order $<$ so that $(i,i') \in E$ and $(i',i) \notin E$  whenever $i \in I_{j}$ and $i' \in I_{j'}$ where $j,j' \in J$ with $j' < j$   
		
		
		
	\end{enumerate}
\end{prop}

\begin{proof}
	Let $b_1 b_2 \cdots b_n$ be alternating product of elements of  $(\mathcal{B}_j)_{j \in J}$ where each $b_r \in \mathcal{B}_{j_r}$ is an alternating product of elements of $(\mathcal{A}_i)_{i \in I_{j_r}}$.   
	%
	Hence, each $b_r$ is of the form  
	%
	\[b_r = a_{(m_0+\cdots+m_{r-1})+1} a_{(m_0+\cdots+m_{r-1})+2}   \cdots  a_{(m_0+\cdots+m_{r-1})+m_r} \]
	%
	with $a_k \in \mathcal{A}_{i_k}$, $i_k \in I_{j_r}$, and $i_{k} \neq i_{k+1}$.   
	%
	Note that we have assumed $m_0=0$.  
	%
	Take $\ell_k = m_0 + m_1 + m_2 + \cdots + m_k$ and $m=\ell_n$. 
	%
	Thus, we have %letting $m = m_1 + \cdots + m_n$, we obtain  
	$$\varphi(b_1 b_2 \cdots b_n) \ = \ \varphi( a_1 a_2 \cdots a_{m} ).$$
	%
	%Let us show how different hypothesis lead to boolean, monotone, and tensor
	%
	%
	%
	
	\noindent \textbf{Proof of \emph{(i)} .} Suppose $(i,i') \in E$ whenever $i \in I_{j}$ and $i' \in I_{j'}$ where $j,j' \in J$ with $j \neq j'$. 
	%
	Put $\bm{j}=(j_1,j_2,\ldots,j_n)$ and $S_U = \cup_{r \in U}[l_{r-1}+1,l_r]$ for each $U \in \ker[\bm{j}] \in P(n)$. 
	%
	%Additionally, we let $j_U$ denote the common value $j_r$ for $r \in U$.  
	%
	We will show that $\pi_U = \ker_G[\bm{i}|_{S_U}]$ is a subset of $\pi = \ker_G[\bm{i}]$ for each $U \in \ker[\bm{j}]$. 
	%
	%
	%
	%
	
	Take any $U \in \ker[\bm{j}]$.
	%
	Suppose $k_1 \sim_{\pi_U} k_2$ with $k_1 , k_2 \in S_U$ with $k_1 < k_2$ and let $\ell \in [n]$ so that $k_1 < \ell < k_2$ and $i_{\ell} \neq i_{k_1}$.  
	%
	Let $r_1,t \in [n]$ with $r_1 \leq t$ such that $k_1 \in [\ell_{r_1-1}+1,\ell_{r_1}]$ and  $\ell \in [\ell_{t-1}+1,\ell_t]$, so we have $i_{k_1} \in I_{j_r}$ and $i_{\ell} \in I_{j_t}$. 
	%
	Note that $\ell \in S_{U}$ if and only if $j_t = j_{r_1}$. 
	%
	Thus, if $\ell \notin S_U$, then $(i_\ell,i_{k_1}) \in E$ by hypothesis since $j_t \neq j_{r_1}$; on the other hand, if $\ell \in S_U$, then $(i_\ell,i_{k_1}) \in E$ since $k_1 \sim_{\pi_U} k_2$. 
	%
	In any case, we obtain $k_1 \sim_{\pi} k_2$. 
	%
	Suppose now $k_1 \in S_U$ and $k_2 \in [m]\setminus S_U$.
	%
	%Note that . 
	%
	Let $r_1, r_2 \in [n]$ such that $k_1 \in [l_{r_1-1}+1,l_{r_1}] $ and $k_2 \in [l_{r_2-1}+1,l_{r_2}]$, so we have $i_{k_1} \in I_{j_{r_1}}$ and $i_{k_2} \in I_{j_{r_2}}$. 
	%
	Since $k_2 \notin S_U$, we have $j_{r_2} \neq j_{r_1}$, and hence $i_{k_1} \neq i_{k_2}$ and $k_1 \nsim_{\pi} k_2$ since $I_{j_{r_1}}$ and $I_{j_{r_2}}$ are disjoint. 
	%
	Thus, Lemma \ref{lem_weak_bm_ver_kernel} implies $\ker_G[\bm{i}|_{S_U}]$ is a subset of $\ker_G[\bm{i}]$. 
	%
	%
	%
	%
	
	Now, since $\ker_G[\bm{i}]$ and $\ker_G[\bm{i}|_{S_U}]$ are partitions of $[m]$ and $S_U$, respectively, and $[m] = \cup_{U \in \ker[\bm{j}]} S_U$, the fact that $\ker_G[\bm{i}|_{S_U}]$ is a subset of $\ker_G[\bm{i}]$ for every $U \in \ker[\bm{j}]$ implies the partition $\ker_G[\bm{i}]$ is the disjoint union of $\ker_G[\bm{i}|_{S_U}]$ with $U \in \ker[\bm{j}]$. 
	%
	It follows from BMT independence that    
	%
	\[
	\varphi( b_1 b_2 \cdots b_n )
	\ \ =  
	\prod_{U \in \ker(\bm{j})} \prod_{ V \in \ker_G(\bm{i}|_{S_U})} \varphi \left( (a_k) | V  \right)
	\ \ =  \prod_{U \in \ker(\bm{j})}  \varphi \left( (b_r) | U  \right). 
	\]
	%
	Therefore,  $(\mathcal{B}_j)_{j \in J}$ are tensor independent. \ \\
	
	
	
	\noindent \textbf{Proof of \emph{ii)}
		.} Suppose now that $(i,i') \notin E$ whenever $i \in I_{j}$ and $i' \in I_{j'}$ with $j,j' \in J$ and $j \neq j'$. 
	%
	Let $(\ell,\ell') = (\ell_0+1,\ell_1+1 )$ and $S = \{ \ell, \ell+1, \ldots, \ell'-1 \}$. 
	%
	Let $\pi$, $\pi_S$, and $\pi_{S^c}$ denote $\ker_G[\bm{i}]$, $\ker_G[\bm{i}|_S]$, and $\ker_G[\bm{i}|_{S^c}]$, respectively.  
	%
	We will show that the partitions $\ker_G[\bm{i}|_S]$ and $\ker_G[\bm{i}|_{S^c}]$ are contained in  $\ker_G[\bm{i}]$. 
	%
	%
	%
	%
	
	Since $S$ and $S^c$ are sub-intervals of $[m]$, if either $k_1 \sim_{\pi_S} k_2$ with $k_1,k_2 \in S$ or $k_1 \sim_{\pi_{S^c}} k_2$ with $k_1,k_2 \in S^c$, then $k_1 \sim_\pi k_2$.
	%
	Supose now $k_1 \in S$ and $k_2 \in S^c$. 
	%
	Thus, we have $i_{k_1} \in I_{j_1}$ and $i_{k_2} \in I_{j_r}$ for some $r \geq 2$. 
	%
	If $r=2$, then $i_{k_1} \neq i_{k_2}$ since $I_{j_1}$ and $I_{j_2}$ are disjoint due to $j_1 \neq j_2$;
	%
	on the other hand, if $r >3$, then either $i_{k_1} \neq i_{k_2}$ or $i_{k_1} = i_{k_2}$ and $k_1 < \ell' \leq \ell_2 < k_2$ with  $(i_{\ell'},i_{k_1}) \notin E$ since $i_{\ell'} \in I_{j_2}$  and $j_1 \neq j_2$. 
	%
	In any case, we obtain $k_1 \nsim_{\pi} k_2$. 
	%
	It follows from Lemma \ref{lem_weak_bm_ver_kernel} that $\ker_G[\bm{i}|_S]$ and $\ker_G[\bm{i}|_{S^c}]$ are subsets of $\ker_G[\bm{i}]$. 
	%
	%
	%
	%
	%
	
	Hence, we obtain $\varphi(b_1 b_2 \cdots b_n)  = \varphi(b_1) \varphi(b_2 \cdots b_n)$ by Lemma \ref{lem_weak_bm}.  
	%	
	Repeating the same argument for $(\ell,\ell')=(\ell_1+1,\ell_2+1),\ldots,(\ell_{n-2}+1,\ell_{n-1}+1)$, we obtain $\varphi(b_1 b_2 \cdots b_n)  = \varphi(b_1) \varphi(b_2) \cdots \varphi(b_n)$. 
	%\begin{enumerate} \item[(B)] \makebox[\linewidth]{$\varphi(b_1 b_2 \cdots b_n)  = \varphi(b_1) \varphi(b_2) \cdots \varphi(b_n).$} \end{enumerate} Therefore, $(\mathcal{B}_j)_{j \in J}$ are boolean independent. 
	Therefore, $(\mathcal{B}_j)_{j \in J}$ are boolean independent. \ \\
	%
	%
	%
	%
	
	
	\noindent \textbf{Proof of \emph{iii)}
		.} Finally, suppose now that $J$ has a total order $<$ so that $(i,i') \in E$ and $(i',i) \notin E$  whenever $i \in I_{j}$ and $i' \in I_{j'}$ where $j,j' \in J$ with $j' < j$. 
	%
	Assume first $ j_1 > \cdots > j_{r-1} > j_r < j_{r+1} < \cdots < j_n$ for some $r$. 
	%
	The same recursive argument used for boolean independence above can be applied in this case, first for $(\ell,\ell') =(\ell_0+1,\ell_1+1),\ldots,(\ell_{r-2}+1,\ell_{r-1}+1)$, to get $\varphi(b_1 b_2 \cdots b_n)  = \varphi(b_1)  \cdots \varphi(b_{r-1}) \varphi(b_r \cdots b_n)$, and second for $(\ell',\ell) =(\ell_{n-1}+1,\ell_n),\ldots,(\ell_{r-1}+1,\ell_{r})$.
	%
	Note that we go in decreasing order from $n$ to $r$ in the latter case. 
	%
	We then obtain $$\varphi(b_1 b_2 \cdots b_n)  = \varphi(b_1)  \cdots \varphi(b_{r-1}) \varphi(b_r)  \cdots \varphi(b_n).$$ 
	%
	%
	%
	%
	
	Assume now $j_{r-1} < j_r > j_{r+1}$ for some $r$. 
	%
	Let $\ell = \ell_{r-1}+1$, $\ell' = \ell_r+1$, and $S = \{ \ell, \ell+1, \ldots, \ell'-1 \}$. 
	%
	We will show that $\ker_G[\bm{i}|_S]$ and $\ker_G[\bm{i}|_{S^c}]$ are subsets of $\ker_G[\bm{i}]$. 
	%
	%
	%
	
	Since $S$ is a sub-interval of $[m]$, we obtain $k_1 \sim_{\pi} k_2$ for any $k_1,k_2 \in S$ with  $k_1 \sim_{\pi_S} k_2$. 
	%
	Suppose now $k_1 \in S$ and $k_2 \in S^c$. 
	%
	Thus, we have $i_{k_1} \in I_{j_r}$ and $i_{k_2} \in I_{j_t}$ for some $ t \in \{1,\ldots,r-1,r+1,\ldots,n\}$.    
	%
	If $t=r+1$, then $i_{k_1} \neq i_{k_2}$ since $I_{j_r}$ and $I_{j_{r+1}}$ are disjoint due to $j_r \neq j_{r+1}$;
	%
	on the other hand, if $t>r+1$, then either $i_{k_1} \neq i_{k_2}$ or $i_{k_1} = i_{k_2}$ and $k_1 < \ell' \leq \ell_{r+1} < k_2$ with  $(i_{\ell'},i_{k_1}) \notin E$ since $i_{k_1} \in I_{j_r}$, $i_{\ell'} \in I_{j_{r+1}}$, and $j_{r+1} < j_{r}$. 
	%
	Similar arguments work if $t \leq r-1$. 
	%
	In any case, we obtain $k_1 \nsim_{\pi} k_2$. 
	%
	It then follows from Lemma \ref{lem_weak_bm_ver_kernel} that $\ker_G[\bm{i}|_S]$ is contained in $\ker_G[\bm{i}]$. 
	%
	%
	%
	%
	
	To obtain that $\ker_G[\bm{i}|_{S^c}]$ is also contained in $\ker_G[\bm{i}]$, it only remains to prove that $k_1 \sim_\pi k_2$ whenever $k_1 \sim_{\pi_{S^c}} k_2$ with $k_1,k_2 \in S^c$. 
	%
	So, let us assume $k_1,k_2 \in S^c$ with $k_1 \sim_{\pi_{S^c}} k_2$. 
	%
	If $k_1 , k_2 \leq \ell-1$ or $\ell' \leq k_1,k_2$, then $k_1 \sim_\pi k_2$ since $\{1,2,\ldots,\ell-1\}$ and $\{\ell',\ell'+1,\ldots,m\}$ are sub-intervals of $[m]$. 
	%
	Thus, without loss of generality, we can assume $k_1 \leq \ell-l$ and $\ell' \leq k_2$.
	%
	The condition $k_1 \sim_{\pi_{S^c}} k_2$ means $i_{k_1} = i_{k_2}$ and $(i_{\ell''},i_{k_1}) \in E$ for $k_1 < \ell'' < k_2$ with $i_{\ell''} \neq i_{k_1}$ and $\ell'' \in S^c$. 
	%
	Hence, to obtain $k_1 \sim_\pi k_2$, we only need to show that $(i_{\ell''},i_{k_1}) \in E$ for $k_1 < \ell'' < k_2$ with $i_{\ell''} \neq i_{k_1}$ and $\ell'' \in S$. 
	%
	Take $1 \leq t \leq r-1$ so that $i_{k_1} \in I_{j_{t}} $. 
	%
	Now, if $i_{k_1} = i_{\ell-1}$ and $\ell'' \in S$, then $j_t = j_{r-1}$ and $i_{\ell''} \in I_{j_{r}}$, and hence $i_{k_1}\neq i_{\ell''}$ since $I_{j_{r-1}}$ and $I_{j_{r}}$ are disjoint and $(i_{\ell''},i_{k_1}) \in E$ since $j_{r-1} < j_r$.
	%
	On the other hand, if $i_{k_1} \neq i_{\ell-1}$, then $k_1 < \ell -1 < \ell' \leq k_2$ with $\ell - 1 \in S^c$ and $i_{\ell-1} \in I_{j_{r-1}}$, so we must have  $(i_{\ell-1},i_{k_1}) \in E$ since  $k_1 \sim_{\pi_{S^c}} k_2$; however, $(i_{\ell-1},i_{k_1}) \in E$  only if $j_t \leq j_{r-1}$, so we obtain $j_t < j_r$, and hence  $(i_{\ell''},i_{k_1}) \in E$  for any $\ell'' \in S$ since $i_{\ell''} \in I_{j_{r}}$ and $i_{k_1} \in I_{j_t}$. 
	%
	In any case, we get $k_1 \sim_\pi k_2$, and therefore $\ker_G[\bm{i}|_{S^c}]$ is a subset of  $\ker_G[\bm{i}]$ by Lemma \ref{lem_weak_bm_ver_kernel}. 
	%
	%
	%
	
	It follows from Lemma \ref{lem_weak_bm} that $\varphi(b_1 \cdots b_n) = \varphi(b_r)  \varphi(b_1 \cdots b_{r-1}b_{r+1}\cdots b_n)$. 
	%
	Therefore,  $(\mathcal{B}_j)_{j \in J}$ are monotone independent. 
	%
\end{proof}

\end{document}
